% Part: first-order-logic
% Chapter: syntax-and-semantics
% Section: subformulas

\documentclass[../../../include/open-logic-section]{subfiles}

\begin{document}

\olfileid{fol}{syn}{sbf}

\olsection{\printtoken{P}{subformula}: formuledelen}

\begin{explain}
Vaak willen we praten over de !!{formula}s waaruit een andere
!!{formula} is opgebouwd. We noemen die haar \emph{!!{subformula}s}:
het zijn formuledelen die zelf ook formules zijn. Iedere !!{formula}
telt als !!a{subformula} van zichzelf. Een subformule van $!A$ die niet
$!A$ zelf is, heet een \emph{echte !!{subformula}}.
\end{explain}

\begin{defn}[Direct !!^{subformula}]
Als $!A$ !!a{formula} is, noemen we de \emph{directe
!!{subformula}s} van $!A$ de formuledelen die één bouwstap lager
liggen. Dit zijn ze:
\begin{enumerate}
\item Atomaire !!{formula}s hebben geen directe !!{subformula}s.

\tagitem{prvNot}{\indcase{!A}{\lnot !B}{De enige directe
    !!{subformula} van $\indfrm$ is~$!B$.}}{}

\item \indcase{!A}{(!B \ast !C)}{De directe !!{subformula}s van
  $\indfrm$ zijn $!B$ en $!C$ ($\ast$ mag elk logisch verbindingsteken
  met twee plaatsen zijn).}

\tagitem{prvAll}{\indcase{!A}{\lforall[x][!B]}{De enige directe
    !!{subformula} van $\indfrm$ is~$!B$.}}{}

\tagitem{prvEx}{\indcase{!A}{\lexists[x][!B]}{De enige directe
    !!{subformula} van $\indfrm$ is~$!B$.}}{}
\end{enumerate}
\end{defn}

\begin{defn}[Echt !!^{subformula}]
Als $!A$ !!a{formula} is, vinden we de \emph{echte !!{subformula}s}
van $!A$ met de volgende regels. De regels gebruiken zichzelf opnieuw:
\begin{enumerate}
\item Atomaire !!{formula}s hebben geen echte !!{subformula}s.

\tagitem{prvNot}{\indcase{!A}{\lnot !B}{De echte !!{subformula}s van
    $\indfrm$ zijn~$!B$ en alle echte !!{subformula}s van~$!B$.}}{}

\item \indcase{!A}{(!B \ast !C)}{De echte !!{subformula}s van
  $\indfrm$ zijn $!B$ en $!C$. Daarbij komen alle echte
  !!{subformula}s van $!B$ en alle echte !!{subformula}s van~$!C$.}

\tagitem{prvAll}{\indcase{!A}{\lforall[x][!B]}{De echte
    !!{subformula}s van $\indfrm$ zijn~$!B$ en alle echte
    !!{subformula}s van~$!B$.}}{}

\tagitem{prvEx}{\indcase{!A}{\lexists[x][!B]}{De echte
    !!{subformula}s van $\indfrm$ zijn~$!B$ en alle echte
    !!{subformula}s van~$!B$.}}{}
\end{enumerate}
\end{defn}

\begin{defn}[!!^{subformula}]
De !!{subformula}s van $!A$ zijn $!A$ zelf en al haar echte
!!{subformula}s.
\end{defn}

\begin{explain}
De twee definities doen niet precies hetzelfde. Bij directe
!!{subformula}s kijken we alleen één bouwstap omlaag. Voor iedere
mogelijke vorm van een formule~$!A$ staat er rechtstreeks welke
directe !!{subformula}s bij~$!A$ horen. Het is dus een lijst van
gevallen. Die gevallen volgen dezelfde bouwregels als de verzameling
!!{formula}s. Bij echte !!{subformula}s volgen we die bouwregels ook,
maar gaan we verder omlaag. Naast de kleinere !!{formula}s $!B$ en
$!C$ nemen we ook alle echte !!{subformula}s die zij zelf hebben. De
definitie gebruikt zichzelf dus opnieuw; daarom heet zij
\emph{recursief}. Algemeen werkt dit zo voor een functie op een
verzameling die stap voor stap is opgebouwd (hier de !!{formula}s): in
een geval van de definitie mag de functie zichzelf gebruiken. Dat mag
alleen voor invoer die al ``eerder'' aan bod kwam. Bij $(!B \ast !C)$
gebruiken we daarom alleen de echte !!{subformula}s van de ``eerdere''
!!{formula}s $!B$ en $!C$.
\end{explain}

\begin{prop}
\ollabel{prop:subfrm-trans}
Stel dat $!B$ een subformule van $!A$ is en $!C$ een subformule van
$!B$. Dan is $!C$ ook een subformule van $!A$. De relatie ``is een
subformule van'' is dus transitief.
\end{prop}

\begin{prob}
Bewijs \olref[fol][syn][sbf]{prop:subfrm-trans}.
\end{prob}

\begin{prop}
\ollabel{prop:count-subfrms}
Stel dat $!A$ een formule met $n$ logische verbindingstekens en tekens
voor ``voor alle'' of ``er bestaat'' is. Dan heeft $!A$ op zijn meest
$2n+1$ subformules.
\end{prop}

\begin{prob}
Bewijs \olref[fol][syn][sbf]{prop:count-subfrms}.
\end{prob}
\end{document}
